\scanpage{16}
% Pencil foliation: 14, upper right; 15, lower left.
3°) Supposons $S=\operatorname{Spec}(\mathbb{C})$, prenons
$F=T_{\mathrm{Betti}}$\deleted{et $T_{\mathrm{Hdg}}$}.
\margin{\deleted{\unsure{bigr. de Hdg}}}
Alors on \deleted{peut prendre} $K=\mathbb{Q}$. D'autre part
$F\otimes_{\mathbb{Q}}\mathbb{C}\simeq T_{\mathrm{Hdg}}$
\deleted{\unsure{$\otimes_{\mathbb{Q}}\mathbb{C}$}},
donc
$G^{\mathrm{Betti}}\otimes_{\mathbb{Q}}\mathbb{C}\simeq G^{\mathrm{Hdg}}$
\deleted{\unsure{$\otimes_{\mathbb{Q}}\mathbb{C}$}}.
Donc dans $G^{\mathrm{Betti}}\otimes_{\mathbb{Q}}\mathbb{C}$ on trouve
\deleted{\unsure{$i$}} $i_1,i_2$
\diagram{\begin{tikzcd}
i_1,i_2:\mathbb{G}_{\mathrm m}\arrow[r]&G^{\mathrm{Betti}}\otimes_{\mathbb{Q}}\mathbb{C}
\end{tikzcd}}
tels que
\[
 i_1(\lambda)i_2(\lambda)=\unsure{\gamma|g}(\lambda).
\]
De plus on a
\[
 \overline{i_1(\lambda)}=i_2(\unsure{\lambda|\overline\lambda})
\]
ce qui exprime que les deux graduations
de Hodge sont conjuguées complexes l'une de
l'autre.

b) Il y a lieu de supposer que
$\operatorname{Ob}C\unsure{\ni|\supset}\mathbb{Q}(\unsure{-1|1})$.
Alors $\mathbb{Q}(\unsure{-1|1})$ \deleted{\unsure{est}} définit
un $G$-module $F(\mathbb{Q}(\unsure{-1|1}))$ de rang $1$,
d'où un hom
\diagram{\begin{tikzcd}
\varepsilon:G\arrow[r]&\mathbb{G}_{\mathrm m,K}
\end{tikzcd}}
Comme $\mathbb{Q}(-1)$ est de \deleted{\unsure{degré}} poids $2$, on aura
\[
 \boxed{\varepsilon(\unsure{\gamma|g}(\lambda))=\lambda^2}
\]

c) Supposons que sur le foncteur $F$ il y ait
une filtration fonctorielle, compatible avec \unsure{gr.}.
Comment l'exprimer ? On considère
\[
 F'=\operatorname{Gr}(F),
\]
